\documentclass{article}
\usepackage{amsmath}
\usepackage{amssymb}
\usepackage{amsthm}
\usepackage{mathtools}
\usepackage{geometry}
\usepackage{hyperref}

\newcommand{\rk}{\operatorname{rank}}
\newcommand{\GL}{\operatorname{GL}}
\newcommand{\R}{\mathbb{R}}
\newcommand{\Mat}{\mathrm{Mat}}
\newcommand{\Hom}{\operatorname{Hom}}

\theoremstyle{definition}
\newtheorem{theorem}{Theorem}[section]
\newtheorem{proposition}[theorem]{Proposition}
\newtheorem{lemma}[theorem]{Lemma}
\newtheorem{corollary}[theorem]{Corollary}
\newtheorem{definition}[theorem]{Definition}
\newtheorem{remark}[theorem]{Remark}

\title{Task 0: Ambient Algebra and Notation}
\date{}

\begin{document}
\maketitle

\section{Goal}

This file fixes the ambient algebra for the second-conjecture block.

The main point is to translate the direct weight-space cyclic equations
\[
W_{<\ell}W_0W_{>\ell}=0,
\qquad
1\le \ell\le L,
\]
into quiver relations, while keeping the end-to-end path
\[
W_L\cdots W_1
\]
nonzero.

\section{The cyclic quiver}

\begin{definition}[The cyclic quiver $C_{L+1}$]
Let $Q=C_{L+1}$ be the oriented cycle on vertices
\[
Q_0=\{0,1,\dots,L\}
\]
with arrows
\[
a_t:t-1\to t
\qquad
(1\le t\le L),
\qquad
a_0:L\to 0.
\]
\end{definition}

\begin{definition}[Weight-space relation paths]
For each $\ell\in\{1,\dots,L\}$, define the path
\[
p_\ell
:=
a_{\ell-1}\cdots a_1\,a_0\,a_L\cdots a_{\ell+1}.
\]
This is the unique oriented path of length $L$ in $Q$ that omits the arrow
$a_\ell$.
\end{definition}

\begin{remark}
Every path $p_\ell$ passes through the closing arrow $a_0$.

The only oriented path of length $L$ that does \emph{not} pass through $a_0$ is
\[
p_0:=a_L\cdots a_1.
\]
This is the path that corresponds to the end-to-end map of the DLN, so it must
survive in the direct weight-space setup.
\end{remark}

\begin{definition}[The weight-space cyclic algebra]
Let $k$ be a field and let $kQ$ be the path algebra of $Q=C_{L+1}$.
Define the two-sided ideal
\[
I_L^{\mathrm{wt}}
:=
\langle p_1,\dots,p_L\rangle
\subset kQ,
\]
and set
\[
A_L^{\mathrm{wt}}
:=
kQ/I_L^{\mathrm{wt}}.
\]
\end{definition}

\section{Dictionary with weight space}

\begin{definition}[The associated quiver representation of a weight tuple]
Fix vector spaces
\[
V_\ell\cong \R^{d_\ell},
\qquad
0\le \ell\le L,
\]
and a distinguished closing map
\[
W_0\in \Hom(V_L,V_0).
\]
Given a weight tuple
\[
W=(W_1,\dots,W_L),
\qquad
W_\ell\in \Hom(V_{\ell-1},V_\ell),
\]
define the quiver representation
\[
\phi(W):=(\phi_0,\phi_1,\dots,\phi_L)
\]
by
\[
\phi_0:=W_0,
\qquad
\phi_\ell:=W_\ell
\quad
(1\le \ell\le L).
\]
\end{definition}

\begin{proposition}[The weight-space equations are the quiver relations]
Let $W=(W_1,\dots,W_L)$ and let $\phi(W)$ be the associated quiver
representation with fixed closing map $W_0$.
Then, for each $\ell\in\{1,\dots,L\}$,
\[
p_\ell(\phi(W))=W_{<\ell}W_0W_{>\ell}.
\]
Hence
\[
W_{<\ell}W_0W_{>\ell}=0
\quad\Longleftrightarrow\quad
p_\ell(\phi(W))=0.
\]
\end{proposition}

\begin{proof}
By definition, $p_\ell$ is the composition of all arrows of the oriented cycle
except $a_\ell$, in the order
\[
a_{\ell-1}\cdots a_1\,a_0\,a_L\cdots a_{\ell+1}.
\]
After substituting
\[
\phi_0=W_0,
\qquad
\phi_t=W_t \ \text{ for } 1\le t\le L,
\]
this becomes exactly
\[
W_{\ell-1}\cdots W_1\,W_0\,W_L\cdots W_{\ell+1}
=
W_{<\ell}W_0W_{>\ell}.
\]
\end{proof}

\begin{remark}[Why the end-to-end path must survive]
Under the same identification,
\[
p_0(\phi(W))
=
a_L\cdots a_1(\phi(W))
=
W_L\cdots W_1
=
\mu(W).
\]
Therefore imposing $p_0=0$ would force $\mu(W)=0$, which is incompatible with
the direct weight-space rank-$r$ sheet for $r>0$.
\end{remark}

\section{Notation to keep fixed later}

The rest of the second-conjecture block should use the following dictionary
without further renaming:

\begin{itemize}
\item the quiver: $Q=C_{L+1}$;
\item the weight-space ideal: $I_L^{\mathrm{wt}}=\langle p_1,\dots,p_L\rangle$;
\item the algebra: $A_L^{\mathrm{wt}}=kQ/I_L^{\mathrm{wt}}$;
\item the surviving length-$L$ path: $p_0=a_L\cdots a_1$;
\item the specialization of the closing arrow:
\[
a_0 \rightsquigarrow W_0=\Sigma_{XY}P_Q.
\]
\end{itemize}

\section{Deliverables}

This task is complete once the file records:

\begin{itemize}
\item the exact quiver $C_{L+1}$;
\item the exact relation ideal $I_L^{\mathrm{wt}}$;
\item the reason why $p_0=a_L\cdots a_1$ survives;
\item the precise translation from the weight-space equations
  $W_{<\ell}W_0W_{>\ell}=0$ to the quiver relations $p_\ell=0$.
\end{itemize}

\end{document}
